\usepackage[utf8]{inputenc}
\usepackage[T1]{fontenc}
\usepackage[spanish,es-nodecimaldot]{babel}
\usepackage[a4paper,margin=2.8cm,headheight=15pt]{geometry}
\usepackage{lmodern}
\usepackage{microtype}
\usepackage{setspace}
\usepackage{graphicx}
\usepackage{booktabs}
\usepackage{tabularx}
\usepackage{longtable}
\usepackage{array}
\usepackage{float}
\usepackage{enumitem}
\usepackage{amsmath,amssymb}
\usepackage{xcolor}
\usepackage{listings}
\usepackage{csquotes}
\usepackage[hidelinks]{hyperref}
\usepackage{bookmark}
\usepackage{fancyhdr}

\onehalfspacing
\setlist{itemsep=0.25em,topsep=0.35em}
\setlength{\parindent}{1.2em}
\setlength{\parskip}{0.25em}

\definecolor{codegray}{gray}{0.95}
\lstset{
  basicstyle=\ttfamily\small,
  backgroundcolor=\color{codegray},
  frame=single,
  breaklines=true,
  columns=fullflexible,
  keepspaces=true
}

\pagestyle{fancy}
\fancyhf{}
\fancyhead[L]{Plataforma para la experimentación con QML}
\fancyhead[R]{\thepage}
\renewcommand{\headrulewidth}{0.4pt}
\renewcommand{\footrulewidth}{0pt}

\newcommand{\pending}[1]{%
  \par\medskip\noindent\textit{Contenido pendiente: #1}\par\medskip
}

\newcommand{\placeholderfigure}[2]{%
  \begin{figure}[H]
    \centering
    \fbox{\parbox[c][#2][c]{0.85\textwidth}{\centering
      \textit{Figura pendiente: #1}}}
    \caption{#1}
  \end{figure}
}

\newcommand{\placeholdertable}[1]{%
  \begin{table}[H]
    \centering
    \fbox{\parbox{0.85\textwidth}{\centering
      \textit{Tabla pendiente: #1}}}
    \caption{#1}
  \end{table}
}
